\documentclass[11pt,a4paper]{article}

\usepackage[utf8]{inputenc}
\usepackage[T1]{fontenc}
\usepackage[spanish]{babel}
\usepackage{amsmath,amssymb}
\usepackage{geometry}
\usepackage{enumitem}
\usepackage{fancyhdr}
\usepackage{titlesec}
\usepackage{array}

\geometry{margin=2.5cm}

\setlength{\headheight}{14pt}
\pagestyle{fancy}
\fancyhf{}
\rhead{Aplicaciones Computacionales en Negocios -- UTDT}
\lhead{Ejercicios adicionales (2.\textsuperscript{o} parcial)}
\cfoot{\thepage}

\titleformat{\section}{\large\bfseries}{\thesection.}{0.6em}{}
\titlespacing*{\section}{0pt}{1.6em}{0.8em}

\newcommand{\ejercicio}[1]{\par\medskip\noindent\textbf{#1}\par\nopagebreak\smallskip}

\setlist[enumerate,1]{label=\alph*),leftmargin=2em,itemsep=3pt,topsep=4pt}
\setlist[itemize,1]{label=$\bullet$,leftmargin=2em,itemsep=3pt,topsep=4pt}

\newcommand{\R}{\mathbb{R}}
\newcommand{\Z}{\mathbb{Z}}
\newcommand{\dem}{\mathrm{dem}}
\newcommand{\capa}{\mathrm{cap}}
\newcommand{\costo}{\mathrm{costo}}

\setlength{\parindent}{0pt}
\setlength{\parskip}{4pt}

\title{\textbf{Ejercicios adicionales para el segundo parcial}\\[0.3em]
\large Aplicaciones Computacionales en Negocios -- UTDT}
\author{}
\date{}

\begin{document}

\maketitle
\thispagestyle{fancy}

\noindent
Estos ejercicios cubren los temas de la segunda parte de la materia: las clases
teóricas 13 a 22, las prácticas 8 a 12 y la guía 4. Están pensados como práctica
adicional a los modelos de parcial, y siguen el mismo formato que el parcial
real: formular modelos de programación lineal entera, modificar un modelo dado
para incorporar una nueva condición, y analizar conceptualmente el efecto de un
cambio en el modelo. No se incluyen respuestas.

\medskip
\noindent\textbf{Temas cubiertos:}
\begin{enumerate}[label=\arabic*.,leftmargin=2em,itemsep=1pt]
  \item Programación lineal: modelado, poliedros y geometría.
  \item Programación lineal entera: relajación lineal, desigualdades válidas,
  cortes y \emph{branch and bound}.
  \item Restricciones de modelado: lógica proposicional, \emph{big-M},
  linealizaciones.
  \item Problemas de \emph{lot sizing}.
  \item Problemas de \emph{staffing} y de asignación.
  \item Problemas de \emph{facility location}.
  \item Diseño de \emph{fixtures} deportivos.
  \item Ejercicios integradores: \emph{scheduling} y modelos sobre grafos.
\end{enumerate}

%=====================================================================
\section{Programación lineal: modelado, poliedros y geometría}
%=====================================================================

\ejercicio{1.1}
Una fábrica de bebidas produce tres variedades de gaseosa a partir de dos
insumos: jarabe y agua gasificada. Cada 1000 litros de la variedad $A$ requiere
120 litros de jarabe y deja un margen de USD 85; la variedad $B$ requiere 95
litros de jarabe y deja USD 70; la variedad $C$ requiere 150 litros de jarabe y
deja USD 100. Disponemos de 9000 litros de jarabe por semana y de una capacidad
de embotellado de 50\,000 litros semanales en total.
\begin{enumerate}
  \item Formular un modelo de programación lineal para maximizar el margen
  semanal.
  \item Agregar la condición de que la producción conjunta de $B$ y $C$ no
  supere el 60\,\% de la producción total. Mostrar explícitamente cómo se
  linealiza esa condición.
  \item Agregar la condición de que, si se produce la variedad $C$, se deba
  pagar un costo fijo semanal de USD 400 por la habilitación de la línea
  correspondiente. ¿Sigue siendo un modelo de programación lineal?
  \item Generalizar el modelo del primer punto a $n$ variedades, dándoles
  nombres adecuados a los conjuntos y los parámetros.
\end{enumerate}

\ejercicio{1.2}
Consideremos el poliedro $P \subseteq \R^2$ definido por las desigualdades
\[
x_1 + 2x_2 \le 8, \qquad 3x_1 + x_2 \le 9, \qquad x_1 \ge 0, \qquad x_2 \ge 0 .
\]
\begin{enumerate}
  \item Graficar $P$ e identificar todos sus vértices.
  \item Resolver gráficamente $\max\{x_1 + 3x_2 : x \in P\}$ y
  $\max\{3x_1 + x_2 : x \in P\}$. ¿En qué caso la solución óptima es única?
  \item Enunciar el teorema de Minkowski y explicar por qué justifica que el
  método Simplex recorra únicamente vértices.
  \item Dar un ejemplo de función objetivo para la cual el problema sobre $P$
  tenga infinitas soluciones óptimas, y otro ejemplo de poliedro (no acotado)
  sobre el cual un problema de maximización no tenga óptimo finito.
\end{enumerate}

\ejercicio{1.3}
Consideremos un grafo dirigido con nodos $\{1,\dots,6\}$, nodo origen $s = 1$,
nodo destino $t = 6$ y arcos con las siguientes capacidades:
$1\!\to\!2$: 8, $1\!\to\!3$: 10, $3\!\to\!2$: 3, $2\!\to\!4$: 2, $2\!\to\!5$: 7,
$3\!\to\!5$: 12, $5\!\to\!4$: 4, $4\!\to\!6$: 10, $5\!\to\!6$: 8.
\begin{enumerate}
  \item Plantear un modelo de programación lineal para calcular el flujo máximo
  entre $s$ y $t$ respetando las capacidades. \emph{Pista:} incluir una variable
  por cada arco.
  \item ¿Hace falta pedir que las variables sean enteras para obtener una
  solución entera? Justificar.
  \item ¿Cómo se modifica el modelo si además cada arco tiene un costo por
  unidad de flujo y queremos enviar exactamente 10 unidades a costo mínimo?
\end{enumerate}

\ejercicio{1.4}
Un avión tiene tres compartimientos de carga, con capacidades de 12, 18 y 10
toneladas y volúmenes máximos de 700, 900 y 500 m$^3$ respectivamente. Tenemos
un conjunto $Q$ de pedidos; de cada pedido $q \in Q$ conocemos su peso total
$w_q$, su volumen total $v_q$ y su \emph{profit} por tonelada $p_q$. Se puede
aceptar cualquier fracción de un pedido, y el peso, el volumen y el
\emph{profit} son proporcionales a la fracción aceptada.
\begin{enumerate}
  \item Formular un modelo de programación lineal para decidir qué fracción de
  cada pedido aceptar y en qué compartimiento ubicarla, maximizando el
  \emph{profit} total.
  \item Agregar restricciones para que el porcentaje utilizado de la capacidad
  de peso sea el mismo en los tres compartimientos (estabilidad del avión).
  \item ¿Cómo cambia el modelo si los pedidos no se pueden fraccionar?
  ¿Sigue siendo un problema de programación lineal?
\end{enumerate}

%=====================================================================
\section{Programación lineal entera: relajación, cortes y branch and bound}
%=====================================================================

\ejercicio{2.1}
Consideremos el problema
\[
\max \{\, x_1 + 2x_2 \;:\; 5x_1 + 6x_2 \le 30,\ -6x_1 + 7x_2 \le 19,\
2x_2 \ge 1,\ x_1, x_2 \in \Z_+ \,\}.
\]
\begin{enumerate}
  \item Graficar la relajación lineal y marcar todos los puntos enteros
  factibles.
  \item Resolver gráficamente la relajación lineal. ¿Es entero el óptimo?
  ¿Qué información aporta el valor óptimo de la relajación sobre el valor
  óptimo del problema entero?
  \item Proponer una desigualdad válida que corte el óptimo de la relajación.
  ¿Es necesario que una desigualdad válida tenga coeficientes enteros?
  \item Explicar por qué redondear el óptimo de la relajación no es, en general,
  una estrategia correcta.
\end{enumerate}

\ejercicio{2.2}
Sea $Q = \{\, x \in \{0,1\}^n : \sum_{i=1}^{n} s_i x_i \le b \,\}$ con
$s_i > 0$ para todo $i$ y $b > 0$.
\begin{enumerate}
  \item Definir qué es un \emph{cubrimiento} (\emph{cover}) y demostrar que, si
  $K$ es un cubrimiento, entonces $\sum_{i \in K} x_i \le |K| - 1$ es una
  desigualdad válida para $Q$.
  \item Mostrar que si $K \subset K'$ son ambos cubrimientos, entonces la
  desigualdad asociada a $K$ implica la asociada a $K'$. ¿Qué consecuencia tiene
  esto a la hora de decidir qué cortes agregar al modelo?
  \item Explicar por qué no conviene agregar de entrada todas las desigualdades
  de cubrimiento. ¿Qué alternativa se usa en la práctica?
  \item Con $n = 5$, $s = (7, 6, 5, 4, 3)$ y $b = 12$, listar todos los
  cubrimientos minimales y escribir sus desigualdades.
\end{enumerate}

\ejercicio{2.3}
Consideremos el modelo de conjunto independiente máximo sobre un grafo
$G = (V,E)$:
\[
\max \sum_{i \in V} x_i \quad \text{sujeto a} \quad x_i + x_j \le 1 \ \
\forall\, ij \in E, \qquad x_i \in \{0,1\} \ \ \forall\, i \in V .
\]
\begin{enumerate}
  \item Mostrar que la relajación lineal de este modelo siempre admite la
  solución $x_i = 1/2$ para todo $i \in V$. ¿Qué cota da esa solución? ¿Qué tan
  mala puede ser esa cota respecto del óptimo entero?
  \item Demostrar que si $K \subseteq V$ es una clique de $G$, entonces
  $\sum_{i \in K} x_i \le 1$ es una desigualdad válida.
  \item Explicar por qué la desigualdad de clique mejora la relajación lineal
  respecto de las restricciones por arista, usando como ejemplo un triángulo.
  \item Escribir un modelo de programación lineal entera para el problema de
  clique máxima, y adaptarlo al caso en que cada vértice $v$ tiene un peso
  $w_v$ y al caso en que cada arista $e$ tiene un peso $w_e$.
\end{enumerate}

\ejercicio{2.4}
Responder brevemente, justificando en cada caso.
\begin{enumerate}
  \item ¿Qué es la relajación lineal de un modelo de programación lineal entera
  y qué tipo de cota proporciona en un problema de minimización?
  \item En el algoritmo de \emph{branch and bound}, ¿qué rol cumplen la cota
  dual y la mejor solución entera hallada hasta el momento ($x^{BEST}$)?
  ¿Cuándo se puede podar un subproblema?
  \item ¿Qué diferencia hay entre una desigualdad válida y un corte?
  \item Dos formulaciones distintas de un mismo problema tienen exactamente las
  mismas soluciones enteras factibles. ¿Pueden tener distinta dificultad
  computacional? Justificar en términos de la cápsula convexa de los puntos
  enteros.
  \item ¿Por qué se dice que, si conociéramos la descripción por desigualdades
  de la cápsula convexa de los puntos enteros factibles, alcanzaría con resolver
  un problema de programación lineal?
\end{enumerate}

%=====================================================================
\section{Restricciones de modelado}
%=====================================================================

\ejercicio{3.1}
Para cada una de las siguientes fórmulas de lógica proposicional, encontrar un
conjunto de restricciones lineales sobre variables binarias cuyas soluciones
factibles se correspondan con los valores que la hacen verdadera. Pueden ser
necesarias variables auxiliares.
\begin{enumerate}[label=\alph*),itemsep=1pt]
  \item $x \veebar y$ (disyunción exclusiva)
  \item $x \Rightarrow \neg y$
  \item $(x \wedge y) \Rightarrow z$
  \item $z \Leftrightarrow (x \vee y)$
  \item $(x_1 \vee \dots \vee x_n) \Rightarrow (y_1 \wedge \dots \wedge y_m)$
  \item ``a lo sumo dos de las variables $x_1,\dots,x_n$ son verdaderas''
  \item ``si $x_1$ es verdadera, entonces al menos tres de
  $x_2,\dots,x_n$ lo son''
\end{enumerate}

\ejercicio{3.2}
Sea $x \in \R_+$ una variable de producción e $y \in \{0,1\}$ una variable que
indica si la línea está encendida.
\begin{enumerate}
  \item Escribir las restricciones que modelan ``$x$ puede ser positiva
  solamente si $y = 1$'', suponiendo que existe una capacidad máxima $C$.
  \item Escribir las restricciones que modelan la variable semicontinua
  $x \in \{0\} \cup [m, C]$.
  \item Si no hubiera capacidad explícita, ¿qué valor de $M$ usarías en
  $x \le M y$ en un modelo de \emph{lot sizing} de $n$ períodos? ¿Por qué
  conviene elegir el $M$ más chico posible?
  \item Modelar ``$x$ debe ser un múltiplo de $L$ unidades, con a lo sumo $k$
  lotes''.
  \item Supongamos que queremos forzar la equivalencia $y = 1 \iff x > 0$.
  ¿Hace falta agregar la implicación recíproca si el objetivo minimiza un costo
  positivo asociado a $y$? Justificar.
\end{enumerate}

\ejercicio{3.3}
Explicar cómo linealizar cada una de las siguientes situaciones, indicando qué
variables auxiliares se agregan y por qué la linealización es correcta.
\begin{enumerate}
  \item Minimizar $\max_{i \in I} \big( \sum_{j} c_{ij} x_{ij} \big)$.
  \item Maximizar $\min_{i \in I} \big( \sum_{j} c_{ij} x_{ij} \big)$.
  \item Minimizar la diferencia entre el mayor y el menor de los valores
  $\sum_j c_{ij} x_{ij}$, para $i \in I$.
  \item Minimizar $\sum_{i \in I} \max\{0, f_i(x)\}$, con $f_i$ lineal.
  \item El producto $z = x \cdot y$ con $x, y \in \{0,1\}$.
  \item El producto $z = x \cdot y$ con $x \in \{0,1\}$ e $y \in [0, U]$
  continua.
\end{enumerate}

\ejercicio{3.4}
Tenemos un conjunto $I$ de proyectos, cada uno con un beneficio $b_i$ y un costo
$c_i$, y un presupuesto $B$. Sea $x_i \in \{0,1\}$ la variable que indica si el
proyecto $i$ se ejecuta. Escribir las restricciones correspondientes a cada una
de las siguientes condiciones:
\begin{enumerate}[itemsep=1pt]
  \item No se pueden ejecutar los cuatro proyectos del conjunto
  $S = \{1,2,3,4\}$ simultáneamente.
  \item El proyecto 7 sólo se puede ejecutar si se ejecutan el 2 y el 3.
  \item Se debe ejecutar exactamente uno de los proyectos 5 y 6.
  \item Si no se ejecuta ningún proyecto del conjunto $S_1$, entonces se debe
  ejecutar al menos uno del conjunto $S_2$.
  \item Si se ejecuta algún proyecto cuyo costo supera los USD 100\,000,
  entonces se deben ejecutar al menos dos proyectos de costo menor a USD
  20\,000.
  \item No se pueden ejecutar más de tres proyectos cuyo beneficio sea inferior
  al promedio de los $b_i$.
\end{enumerate}

%=====================================================================
\section{Problemas de lot sizing}
%=====================================================================

\ejercicio{4.1}
Consideremos el modelo clásico de \emph{lot sizing} con un único producto:
\[
\min \sum_{t \in T} \big( c_t x_t + a_t s_t \big), \qquad
s_t = s_{t-1} + x_t - d_t \ \ \forall\, t \in T, \qquad
s_0 = 0, \qquad x_t, s_t \ge 0 \ \ \forall\, t \in T .
\]
\begin{enumerate}
  \item ¿Cómo se modifica el modelo si estamos obligados a terminar el período
  $n$ con al menos 150 unidades en stock?
  \item ¿Cómo se modifica el modelo si tenemos un costo fijo $f_t$ por tener
  stock en el depósito en el período $t$? Es decir, pagamos $f_t$ si $s_t > 0$,
  y no lo pagamos si $s_t = 0$.
  \item Tenemos un stock de seguridad $m$ en cada período, que se puede
  incumplir pagando una penalidad de $\Pi$ pesos por período incumplido.
  ¿Cómo se modifica el modelo? ¿Y si la penalidad fuera proporcional a las
  unidades faltantes?
  \item ¿Cómo se modifica el modelo si se permite atrasar el cumplimiento de la
  demanda (\emph{backlogging}), pagando un costo $b_t$ por unidad atrasada en el
  período $t$? Justificar por qué, en el óptimo, las dos variables que
  reemplazan a $s_t$ no pueden ser simultáneamente positivas.
  \item ¿Cómo se modifica el modelo si los costos deben expresarse en valor
  presente usando una tasa de descuento $r$?
\end{enumerate}

\ejercicio{4.2}
Somos fabricantes de cajas de cartón corrugado y nuestro principal insumo es el
papel. Tenemos tres proveedores (A, B y C), un conjunto $T = \{1,\dots,n\}$ de
períodos, la demanda $d_t$ de papel en cada período y el precio por tonelada de
cada proveedor en cada período. Se puede guardar stock de un período al
siguiente.
\begin{enumerate}
  \item Plantear un modelo de programación lineal para determinar cuánto
  comprarle a cada proveedor en cada período, minimizando el costo total.
  \item Si no se pudiera guardar stock, ¿tiene sentido plantear un modelo de
  programación lineal para este problema? Justificar.
  \item Agregar al modelo una capacidad máxima de almacenamiento.
  \item Agregar al modelo la condición de que sólo le podamos comprar al
  proveedor A en a lo sumo cinco períodos, elegidos libremente.
  \item Agregar al modelo la condición de que en cada período le podamos comprar
  a lo sumo a un proveedor.
  \item Cada proveedor exige un pedido mínimo de 5 toneladas por compra.
  ¿Cómo se agrega esta condición?
\end{enumerate}

\ejercicio{4.3}
Un laboratorio fabrica un granel que luego se envasa en blísters, y los blísters
se agrupan en dos presentaciones finales. Cada blíster consume 10 unidades de
granel, la presentación ``$\times 10$'' consume 10 blísters y la presentación
``$\times 20$'' consume 20 blísters. Tenemos $T = \{1,\dots,18\}$ meses y
demanda mensual estimada de cada uno de los cuatro elementos (granel, blíster y
las dos presentaciones finales).
\begin{enumerate}
  \item Formular un modelo de programación lineal entera mixta que determine
  cuánto fabricar de cada elemento en cada mes, respetando la estructura del
  producto y cumpliendo la demanda.
  \item Agregar la condición de que el granel se fabrique en lotes exactos de
  500 unidades y los blísters en lotes de al menos 1000 unidades.
  \item Agregar la condición de que entre dos meses con producción de granel
  deban pasar al menos tres meses.
  \item Agregar la condición de que el granel fabricado en un mes se debe
  blistear en ese mismo mes (es decir, no puede quedar stock de granel).
  \item Cada elemento se fabrica en una máquina que tiene un pool de $H_{kt}$
  horas disponibles en el mes $t$, y la unidad del elemento $i$ consume
  $\tau_i$ horas. Agregar esta restricción.
\end{enumerate}

\ejercicio{4.4}
Una empresa planifica la producción de $P$ productos a lo largo de
$T = \{1,\dots,12\}$ meses, con demanda $d_{it}$, costo unitario $c_{it}$ y
costo de almacenamiento $h_i$.
\begin{enumerate}
  \item Formular el modelo multi-ítem básico.
  \item Agregar una capacidad conjunta de producción de 300 unidades por mes y
  una capacidad conjunta de depósito de $S_{\max}$ unidades.
  \item Se puede tercerizar producción a un costo $g_i$ por unidad, con un tope
  conjunto de 200 unidades por mes y un pedido mínimo de 20 unidades por
  producto tercerizado. Agregar esta posibilidad.
  \item ¿Cuál es la mínima capacidad de depósito que hace factible el problema?
  Explicar cómo se puede obtener ese valor con un único modelo, sin hacer una
  búsqueda manual sobre el valor de $S_{\max}$.
  \item La solución óptima termina el horizonte con stock nulo de todos los
  productos. Explicar por qué ocurre y proponer dos maneras distintas de
  corregir este efecto de horizonte.
\end{enumerate}

%=====================================================================
\section{Problemas de staffing y de asignación}
%=====================================================================

\ejercicio{5.1}
Un banco abre de lunes a viernes de 9 a 18, y conoce la cantidad de cajeros
$r_h$ que necesita en cada hora $h \in \{9,\dots,17\}$ (donde $h$ indica la hora
que comienza en $h$ y termina en $h+1$). Hay dos tipos de cajeros: los
\emph{full-time} trabajan las nueve horas excepto una hora de almuerzo, que el
banco asigna entre las 12 y las 13 o entre las 13 y las 14, y cobran \$800 por
hora (incluida la del almuerzo). Los \emph{part-time} trabajan tres horas
consecutivas y cobran \$500 por hora. El banco no quiere contratar más de cinco
cajeros \emph{part-time}.
\begin{enumerate}
  \item Formular un modelo de programación lineal entera que cubra los
  requerimientos a costo mínimo. Escribir explícitamente la expresión de la
  cantidad de \emph{part-time} activos en la hora $h$.
  \item ¿Cómo se modifica el modelo si cada \emph{part-time} no puede empezar
  antes de las 11?
  \item ¿Cómo se modifica el modelo si la cantidad de \emph{full-time} que
  almuerzan en cada franja no puede diferir en más de dos?
  \item ¿Qué pasaría si relajáramos la integralidad de las variables? ¿Tendría
  sentido la solución obtenida?
\end{enumerate}

\ejercicio{5.2}
Tenemos que conformar un equipo de trabajo a partir de un conjunto $P$ de
personas. Para cada persona $i \in P$ conocemos su sueldo $s_i$ y el resultado
$q_i$ de un test cognitivo. Debemos contratar exactamente 11 personas sin
superar un presupuesto de USD 50\,000 mensuales, maximizando la suma de los
resultados de los tests.
\begin{enumerate}
  \item Formular el modelo. Verificar que cumple las tres reglas de
  presentación de modelos vistas en clase: índices ligados a un cuantificador,
  dominio de las variables especificado, y datos distinguidos de las variables
  de decisión.
  \item Agregar un subconjunto $A \subseteq P$ de personas de contratación
  obligatoria y un conjunto $I \subseteq P \times P$ de pares incompatibles.
  \item Si se elimina la restricción de presupuesto y se busca maximizar la
  cantidad de personas contratadas, ¿qué problema clásico se obtiene?
  \item ¿Qué otros problemas de negocio se pueden representar con el modelo del
  primer punto?
\end{enumerate}

\ejercicio{5.3}
Tenemos un conjunto $P$ de personas y un conjunto $R$ de roles. Para cada
persona $i$ y cada rol $r$ conocemos si está habilitada ($H_{ir} \in \{0,1\}$) y
el puntaje $q_{ir}$ que obtendría. Se requieren exactamente $D_r$ personas en el
rol $r$, cada persona ocupa a lo sumo un rol, y hay un presupuesto $B$ sobre la
suma de los sueldos $s_i$.
\begin{enumerate}
  \item Formular el modelo maximizando el puntaje total.
  \item Agregar: si la persona $u$ es seleccionada, entonces la persona $v$
  también debe serlo.
  \item Agregar: a lo sumo una persona con menos de dos años de experiencia
  (dato $\mathrm{exp}_i$), y si entra alguna persona con un año de experiencia,
  entonces deben entrar al menos dos con tres o más años.
  \item Tres personas $u$, $v$ y $w$ son incompatibles de a pares. Mostrar que
  la relajación lineal admite $x_u = x_v = x_w = 1/2$ y proponer una desigualdad
  válida que elimine esa solución.
  \item ¿Cómo se modifica el modelo si el puntaje de una persona no depende del
  rol que ocupa? ¿Alguna restricción se vuelve redundante?
\end{enumerate}

\ejercicio{5.4}
Estamos a cargo de las operaciones de una aerolínea con base en Aeroparque.
Tenemos un conjunto $V$ de paquetes de vuelos (cada uno comienza y termina en
Aeroparque en un mismo día, y tiene un avión asignado) y un conjunto $P$ de
pilotos. La matriz $M \in \{0,1\}^{|P| \times |V|}$ especifica qué pilotos pueden
realizar cada paquete.
\begin{enumerate}
  \item Formular un modelo de programación lineal entera que maximice la
  cantidad de paquetes con pilotos asignados, sabiendo que cada paquete requiere
  dos pilotos compatibles (aunque es aceptable no asignarlos si no es posible),
  que un piloto no puede trabajar más de cinco días seguidos, que debe tener al
  menos dos días de descanso cada siete días y al menos cinco días de descanso
  cada diez días.
  \item ¿Cómo se modifica el modelo si queremos además equilibrar la carga de
  trabajo, minimizando la máxima cantidad de paquetes asignados a un mismo
  piloto?
  \item ¿Cómo se modifica el modelo si cada paquete tiene una prioridad
  $\rho_v$ y queremos maximizar la prioridad total cubierta?
\end{enumerate}

%=====================================================================
\section{Problemas de facility location}
%=====================================================================

\ejercicio{6.1}
Consideremos el modelo de \emph{facility location} con un conjunto $C$ de
clientes y un conjunto $F$ de \emph{facilities}:
\[
\min \sum_{i \in C}\sum_{j \in F} \costo_{ij}\, x_{ij}
     + \sum_{j \in F} \beta\, y_j, \qquad
\sum_{j \in F} x_{ij} = 1 \ \ \forall\, i \in C, \qquad
x_{ij} \le y_j \ \ \forall\, i \in C,\ \forall\, j \in F,
\]
con $x_{ij}, y_j \in \{0,1\}$ y $\beta > 0$.
\begin{enumerate}[itemsep=1pt]
  \item ¿Qué sucede si se elimina la restricción $x_{ij} \le y_j$?
  \item La formulación permite que $y_j = 1$ aunque ningún cliente use el
  \emph{facility} $j$. ¿Es esto un problema? ¿Cambia la respuesta si $\beta$
  fuera negativo?
  \item Necesitamos abrir por lo menos cinco \emph{facilities} y que cada uno de
  ellos tenga al menos dos clientes asignados. Agregar esta condición.
  \item Los clientes 1, 2 y 3 deben usar tres \emph{facilities} distintos.
  Agregar esta condición.
  \item Conocemos el volumen de compras $v_i$ de cada cliente. Agregar la
  condición de que el volumen total asignado a cada \emph{facility} no supere un
  umbral $\gamma$, y la de que no se abra un \emph{facility} si el volumen
  asignado es menor a $\delta$.
  \item ¿Cómo se modifica el modelo si queremos minimizar el volumen de ventas
  del \emph{facility} con mayor volumen?
  \item ¿Cómo se modifica el modelo si cada \emph{facility} $j$ tiene su propio
  costo de apertura $\beta_j$ y un tope global de presupuesto $B$ para las
  aperturas?
\end{enumerate}

\ejercicio{6.2}
Una empresa produce en una única planta y abastece a 20 clientes mayoristas a
través de distribuidoras regionales. Hay 6 ubicaciones candidatas. Para cada
candidata $i$ conocemos la distancia $D_i$ desde la planta, el costo fijo
mensual $f_i$ de alquiler y, para cada cliente $j$, la distancia $d_{ij}$ desde
la distribuidora $i$. Por cada cliente asignado a una distribuidora se genera un
flujo de abastecimiento desde la planta. Ninguna distribuidora puede atender a
más de 15 clientes, y el presupuesto mensual es $B$.
\begin{enumerate}
  \item Formular un modelo de programación lineal entera que minimice los
  kilómetros totales recorridos (tramo primario más secundario).
  \item Observar que la restricción $\sum_j x_{ij} \le 15\, y_i$ cumple
  simultáneamente dos funciones. ¿Cuáles? ¿Qué pasaría con la relajación lineal
  si en su lugar usáramos $x_{ij} \le y_i$ junto con $\sum_j x_{ij} \le 15$?
  \item ¿Cómo se modifica el modelo si el costo fijo debe aparecer en la función
  objetivo en lugar de en una restricción de presupuesto? ¿En qué unidades
  quedaría la función objetivo y qué problema trae eso?
  \item Agregar la condición de que no se abran dos distribuidoras ubicadas a
  menos de 3 km entre sí, sabiendo que conocemos la distancia $\delta_{ii'}$
  entre cada par de candidatas. ¿Por qué puede ser relevante esta condición?
\end{enumerate}

\ejercicio{6.3}
Trabajamos para una empresa de pagos que recibe $M$ transacciones a procesar en
la nube, y dispone de $N$ servidores. Cada transacción $i$ tiene un tiempo de
procesamiento $t_i$, y cada servidor $j$ tiene un costo fijo de activación $s_j$
y un costo variable $c_j$ por unidad de tiempo.
\begin{enumerate}
  \item Formular un modelo de programación lineal entera que minimice el costo
  total de procesamiento.
  \item ¿Cómo va a ser la solución óptima de esa formulación? ¿Es razonable?
  \item Las transacciones llegan en \emph{batches} y se dispone de un tiempo
  máximo $T$ para procesar un \emph{batch} completo. Modificar el modelo.
  \item ¿Cómo se modifica el modelo si además queremos equilibrar la carga,
  minimizando el tiempo de procesamiento del servidor más cargado?
\end{enumerate}

\ejercicio{6.4}
Debemos ubicar ambulancias en una ciudad. Tenemos un conjunto $L$ de bases
candidatas y un conjunto $C$ de zonas, con población $p_i$ para cada
$i \in C$, y la distancia $d_{ij}$ entre la zona $i$ y la base $j$. Decimos que
la base $j$ \emph{cubre} a la zona $i$ si $d_{ij} \le U$.
\begin{enumerate}
  \item Formular un modelo de programación lineal entera para determinar el
  mínimo número de bases a abrir de modo que todas las zonas queden cubiertas.
  \item Formular un modelo que, abriendo a lo sumo $M$ bases, maximice la
  población cubierta.
  \item ¿Puede el primer modelo ser infactible? ¿Qué se hace en ese caso?
  \item ¿Cómo se modifica el segundo modelo si queremos además que ninguna zona
  quede a más de $2U$ de la base más cercana?
\end{enumerate}

%=====================================================================
\section{Diseño de fixtures deportivos}
%=====================================================================

\ejercicio{7.1}
Consideremos el modelo de \emph{fixture} con $E = \{1,\dots,n\}$ equipos ($n$
par), $F = \{1,\dots,2n-2\}$ fechas y variables binarias $x_{ijf}$ que indican
si el equipo $i$ juega de local contra el equipo $j$ en la fecha $f$, con las
restricciones
\[
\sum_{f \in F} x_{ijf} = 1 \ \ \forall\, i \ne j, \qquad
\sum_{j \in E} \big( x_{ijf} + x_{jif} \big) = 1 \ \ \forall\, i \in E,\
\forall\, f \in F .
\]
\begin{enumerate}[itemsep=1pt]
  \item Escribir la expresión que indica si el equipo $i$ juega de local en la
  fecha $f$, y la que indica si enfrenta al equipo $j$ en la fecha $f$ sin
  importar la localía.
  \item Agregar una restricción para que el equipo 3 tenga un patrón perfecto.
  \item Agregar una restricción para que el equipo 1 no juegue más de tres
  partidos consecutivos de visitante.
  \item Sea $EF = \{1,\dots,5\}$ el conjunto de equipos fuertes. Agregar una
  restricción para que ningún equipo juegue en fechas consecutivas contra dos
  equipos fuertes.
  \item Sea $I$ el conjunto de equipos que juegan copas internacionales y, para
  cada $i \in I$, sea $L_i$ el conjunto de equipos que están a más de una hora
  de avión de $i$. Agregar una restricción para evitar que $i$ juegue dos
  partidos consecutivos de visitante contra equipos de $L_i$ en las primeras 12
  fechas.
  \item Supongamos que agregamos la función objetivo
  $\min \sum_{j \in E} \sum_{f \in F} (x_{1jf} + x_{j1f})$.
  ¿Qué expresa? ¿Tiene sentido para este modelo?
  \item Agregar una función objetivo que maximice la cantidad de partidos
  atractivos (conjunto $A \subseteq E \times E$) programados en las últimas
  cuatro fechas, y otra que los pondere dándole más peso a las fechas más
  tardías.
\end{enumerate}

\ejercicio{7.2}
Una federación organiza un torneo cerrado de ajedrez con un número par $N \ge 6$
de jugadores, formato todos contra todos a una sola rueda ($N-1$ fechas). La
variable binaria $x_{ijt}$ indica que el jugador $i$ juega con blancas contra
$j$ en la fecha $t$.
\begin{enumerate}
  \item Escribir las restricciones de todos contra todos y de un partido por
  jugador por fecha. Observar la diferencia con el modelo de \emph{double
  round-robin}.
  \item Agregar: cada jugador debe disputar la misma cantidad de partidas con
  blancas y con negras, con una diferencia máxima de una partida.
  \item Agregar: ningún jugador puede jugar más de tres partidas consecutivas
  con negras.
  \item Agregar: dos jugadores destacados sólo pueden enfrentarse en las dos
  últimas fechas.
  \item Cada encuentro tiene un puntaje de calidad esperada $w_{ij} \ge 0$.
  Escribir una función objetivo que maximice la calidad de las partidas
  disputadas en dos fechas determinadas $T_{fd} \subset T$.
\end{enumerate}

\ejercicio{7.3}
Sobre \emph{breaks} y patrones de localía.
\begin{enumerate}
  \item Definir qué es un \emph{break} y escribir las variables y restricciones
  que permiten minimizar la cantidad total de \emph{breaks} de un
  \emph{fixture}.
  \item Explicar por qué, dada la función objetivo elegida, alcanza con una sola
  desigualdad por variable de \emph{break}, y por qué esas variables pueden
  declararse continuas.
  \item Demostrar que a lo sumo dos equipos pueden tener patrón perfecto.
  \item Sea $P$ un conjunto de patrones de localía y sea $z_{ip}$ la variable
  binaria que indica si el equipo $i$ usa el patrón $p$. Escribir las
  restricciones que asignan exactamente un patrón a cada equipo y las que ligan
  los patrones con las variables $x_{ijf}$, usando el conjunto $H(f)$ de
  patrones que juegan de local en la fecha $f$.
  \item Explicar la estrategia de resolución en etapas (generación de patrones,
  selección de patrones, resolución del \emph{fixture}) y qué se hace cuando la
  última etapa resulta infactible. Escribir el corte que se agrega al modelo de
  selección a partir de un subconjunto minimal no factible $S$ de asignaciones
  patrón--equipo, y comparar su forma con la de una desigualdad de cubrimiento.
\end{enumerate}

\ejercicio{7.4}
Las eliminatorias sudamericanas para un mundial las juegan 10 equipos, todos
contra todos ida y vuelta, con 3 puntos por partido ganado y 1 por empate. Los
seis primeros clasifican.
\begin{enumerate}
  \item Diseñar un modelo de programación lineal entera que, asignando todos los
  resultados posibles antes de jugada la primera fecha, determine la máxima
  cantidad de puntos con la que el equipo 1 puede \emph{no} clasificar.
  \item Explicar por qué la respuesta a la pregunta ``¿con cuántos puntos tiene
  el equipo 1 garantizada su clasificación?'' es el valor óptimo más uno.
  \item ¿Cómo se modifica el modelo si se desempata por diferencia de gol y
  queremos una garantía que no dependa del desempate?
  \item ¿Hace falta modelar el \emph{fixture} (qué partido se juega en qué
  fecha) para responder esta pregunta? Justificar.
\end{enumerate}

%=====================================================================
\section{Ejercicios integradores}
%=====================================================================

\ejercicio{8.1}
Tenemos un conjunto $T$ de tareas a programar en una máquina que trabaja en
forma continua y procesa una tarea a la vez, sin interrupciones. Conocemos la
duración $d_i$ y la fecha de entrega $e_i$ de cada tarea.
\begin{enumerate}
  \item Escribir un modelo de programación lineal entera mixta que minimice la
  suma de las demoras. \emph{Pista:} usar una variable continua $p_i$ de inicio
  y una binaria $y_{ij}$ para el orden relativo de cada par de tareas.
  \item ¿Son necesarias las variables $y$? ¿Y si el objetivo fuera minimizar el
  \emph{makespan}?
  \item ¿Es necesario resolver este problema con programación lineal entera, o
  existe un algoritmo polinomial? Justificar.
  \item ¿Cómo se modifica el modelo si cada tarea tiene además una fecha de
  inicio más temprana posible?
\end{enumerate}

\ejercicio{8.2}
Consideremos el mismo problema que en el ejercicio anterior, pero con un
conjunto $M$ de $k > 1$ máquinas idénticas.
\begin{enumerate}
  \item Escribir un modelo de programación lineal entera mixta.
  \emph{Pista:} introducir una variable binaria $z_{iu}$ por cada tarea $i$ y
  cada máquina $u$.
  \item ¿Cómo se modifica el modelo si la duración depende de la máquina?
  \item ¿Cómo se modifica el modelo si tenemos un digrafo $G = (T, A)$ de
  precedencias entre las tareas?
  \item ¿Qué valor de \emph{big-M} conviene usar en las restricciones de no
  superposición? Relacionar la elección con la calidad de la relajación lineal.
\end{enumerate}

\ejercicio{8.3}
Dado un grafo $G = (V, E)$:
\begin{enumerate}
  \item Escribir un modelo de programación lineal entera para el problema de
  conjunto dominante mínimo (un conjunto $D \subseteq V$ tal que todo vértice
  está en $D$ o es vecino de algún vértice de $D$).
  \item ¿Cómo cambia el modelo si además pedimos que $D$ sea una clique?
  ¿Siempre tiene solución ese problema?
  \item Escribir un modelo para el problema de coloreo de $G$ con la mínima
  cantidad de colores. \emph{Pista:} usar variables $x_{vk}$ que indiquen si el
  vértice $v$ recibe el color $k$, y variables $w_k$ que indiquen si el color
  $k$ se usa.
  \item El modelo del punto anterior tiene muchas soluciones equivalentes
  (permutaciones de los colores). ¿Por qué es un problema para el
  \emph{solver}? Proponer restricciones que rompan esa simetría.
\end{enumerate}

\ejercicio{8.4}
Consideremos nuevamente la aerolínea del ejercicio 5.4, pero ahora para asignar
aviones a paquetes de vuelos. Tenemos un conjunto $V$ de paquetes y un conjunto
$A$ de aviones. Cada paquete está compuesto por una secuencia de vuelos, con la
cantidad de asientos vendidos en cada uno, y cada avión tiene una cantidad de
asientos disponibles. Si el avión asignado a un paquete tiene menos asientos que
los vendidos en alguno de sus vuelos, se debe hacer un \emph{bump} de los
pasajeros sobrantes.
\begin{enumerate}
  \item Escribir un modelo de programación lineal entera para minimizar la
  cantidad de \emph{bumped passengers}, sabiendo que cada paquete recibe cero o
  un avión (si no recibe ninguno, todos sus pasajeros son \emph{bumped}), que
  dos paquetes que se superponen en el tiempo no pueden recibir el mismo avión,
  y que cada avión debe tener un día de descanso cada 45 días.
  \item ¿Cómo se calcula, en función de los datos, la cantidad de
  \emph{bumped passengers} de un paquete al que se le asigna un avión dado?
  ¿Hace falta una variable para eso?
  \item ¿Cómo se modifica el modelo si el costo de un \emph{bump} depende del
  vuelo?
\end{enumerate}

\ejercicio{8.5}
Preguntas conceptuales de integración. Responder brevemente y con justificación.
\begin{enumerate}[itemsep=1pt]
  \item En un modelo de \emph{facility location} con costo de apertura
  $\beta > 0$, ¿qué sucede con las soluciones óptimas si $\beta$ pasa a ser
  negativo?
  \item ¿Por qué el problema de \emph{lot sizing} sin costos fijos se puede
  resolver con programación lineal, y por qué agregar costos fijos cambia
  sustancialmente la dificultad del problema?
  \item En un modelo de \emph{fixture} sin función objetivo, ¿qué tipo de
  problema estamos resolviendo? ¿Qué implicancias tiene esto para la estrategia
  de resolución?
  \item ¿Qué significa ``fijar características de la solución en lugar de fijar
  decisiones concretas''? Dar un ejemplo de cada tema de la materia.
  \item Mencionar tres estrategias distintas para atacar en la práctica un
  problema NP-completo, y explicar en cuál de ellas se encuadra el uso de un
  \emph{solver} de programación lineal entera.
\end{enumerate}

\end{document}
